\section{Introduction}
\label{sec:introduction}

The concept of \emph{Bechinah} (בחינה, plural: \emph{Bechinot}) occupies a central role in classical Jewish philosophy and Kabbalistic literature, where it denotes an aspect, perspective, or relational modality of revelation. Historically, analysis of \emph{Bechinot} has often been conducted using natural language or metaphysical narrative.

In this work, we present a formal mathematical and computational configuration model titled \textbf{Sod Ha-Bechinah} (סוד הבחינה). Our primary objective is to translate relational intuition into precise mathematical objects, sets, functions, and category-theoretic maps.

We establish three core methodological principles:
\begin{enumerate}
    \item \textbf{Configuration over Isolation}: An aspect ($B$) cannot be modeled in isolation; it emerges strictly from the interaction of Value ($V$), Relation ($R$), and Scale ($S$).
    \item \textbf{Layered Resolution}: Information flows across four distinct resolution layers inspired by \emph{Ta'amim}, \emph{Nekudot}, \emph{Tagin}, and \emph{Otiyot}.
    \item \textbf{Explicit Scope Separation}: Mathematical definitions are strictly distinguished from historical/textual claims and philosophical interpretations.
\end{enumerate}

By doing so, we construct a formal modeling language suitable for computational AI architecture specification.
